\documentclass[11pt,a4paper]{article}
\usepackage{fontspec,amsmath,amssymb,amsthm,microtype}
\setmainfont{lmroman10-regular.otf}[BoldFont=lmroman10-bold.otf,ItalicFont=lmroman10-italic.otf,BoldItalicFont=lmroman10-bolditalic.otf]
\usepackage{xcolor}
\usepackage[most]{tcolorbox}
\usepackage{mathrsfs}
\providecommand{\lt}{<}
\providecommand{\gt}{>}
\setmonofont{DejaVuSansMono.ttf}[BoldFont=DejaVuSansMono-Bold.ttf,ItalicFont=DejaVuSansMono-Oblique.ttf,BoldItalicFont=DejaVuSansMono-BoldOblique.ttf,Scale=0.86]
\newfontfamily\leanfallback{STIXTwoMath-Regular.otf}[Scale=0.86]
\definecolor{leanrule}{gray}{0.62}
\definecolor{leanback}{gray}{0.962}
\definecolor{leantitle}{gray}{0.22}
\definecolor{leanmuted}{gray}{0.38}
\newtcolorbox{leanbox}[1]{enhanced,breakable,colback=leanback,colframe=leanrule,
  boxrule=0.4pt,arc=0.6mm,left=2.4mm,right=2.4mm,top=1.8mm,bottom=1.8mm,
  before skip=0.8em,after skip=1em,fonttitle=\small\bfseries,coltitle=leantitle,
  colbacktitle=leanback,titlerule=0.3pt,title={#1}}
\newlength{\leancharwidth}
\newenvironment{leancode}{\par\smallskip\begingroup\ttfamily\footnotesize\raggedright
  \settowidth{\leancharwidth}{M}\setlength{\parskip}{0pt}\setlength{\parindent}{0pt}
  \frenchspacing}{\par\endgroup\smallskip}
\newcommand{\leanline}[2]{\hangindent=\dimexpr#1\leancharwidth+4\leancharwidth\relax
  \hangafter=1\noindent\hspace*{#1\leancharwidth}#2\par}
\newcommand{\leansource}[1]{\par{\footnotesize\color{leanmuted}#1}\par}
\newcommand{\leanlabel}[1]{\par\smallskip\noindent{\small\bfseries #1}\par}
\usepackage[margin=25mm]{geometry}
\usepackage[unicode,hidelinks]{hyperref}
\hypersetup{pdfauthor={LangevinLean Project}}
\setlength{\emergencystretch}{3em}
\setlength{\parskip}{0.4em}
\setlength{\parindent}{0pt}
\allowdisplaybreaks[2]
\raggedbottom
\title{\texorpdfstring{Weighted weak consistency of the Euler chain}{Weighted weak consistency of the Euler chain}}
\author{LangevinLean Project}
\date{}
\begin{document}\maketitle
Draft mathematical exposition; these notes may contain mistakes.


This section proves a weak consistency estimate for the Euler (unadjusted Langevin) chain that is the discrete input of the Dynkin identity used in the proof of Theorem 2. Let \(H\) be a finite-dimensional real inner product space (in the application \(H=\mathbb R^p\)), \(\gamma\) the standard Gaussian measure on \(H\), \(U:H\to\mathbb R\) a potential with Lipschitz gradient and


\[X_{k+1}=X_k-h\nabla U(X_k)+\sqrt{2h}\,W_k\qquad(k\ge0),\]


where \(W_0,W_1,\dots\) are independent with law \(\gamma\) and independent of \((w,X_0)\) for a bounded random weight \(w\). For a test function \(\varphi\in C^2\) with \(\varphi\) and \(\mathcal L\varphi=\Delta\varphi-\langle\nabla U,\nabla\varphi\rangle\) bounded, write \(T(F,X)=\mathbb E[w\,F(X)]\). Assuming a uniform Gaussian Taylor estimate for \(\varphi\) with constants \(\varepsilon,K\) (proved elsewhere for compactly supported \(\varphi\)), we show


\[\Bigl|T(\varphi,X_n)-T(\varphi,X_0)-\sum_{k=0}^{n-1}h\,T(\mathcal L\varphi,X_k)\Bigr|\le A\,(nh)\,(\varepsilon+K\sqrt h)\,(1+R^2)\]


whenever \(|w|\le A\) almost surely, \(0\le h\le1\), and \(\|\nabla U(X_k)\|_{L_2}\le R\) for \(k\le n\). The weight \(w\) carries arbitrary bounded information about the initial state; it is needed because the Dynkin identity is applied to test functions multiplied by a function of the initial coupling. The proof combines the independence of each new Gaussian increment from the past (including the weight), Fubini's theorem, and a telescoping sum.



Earlier results used here: \cite[Section~3.42]{LangevinLean2026ULASC}.

\section{\texorpdfstring{One Euler step with a fresh Gaussian increment and a bounded weight}{One Euler step with a fresh Gaussian increment and a bounded weight}}


Throughout, \((\Omega,\mathcal F,P)\) is a probability space and \(\gamma\) is the standard Gaussian measure on \(H\).




\textbf{Claim.} (1) Let \(w,F:\Omega\to\mathbb R\) be measurable, \(A\ge0\), \(C\in\mathbb R\), with \(|w|\le A\) almost surely and \(|F|\le C\) everywhere. Then \(wF\) is \(P\)-integrable.




(2) Let \(\varphi:H\to\mathbb R\) be continuous with \(|\varphi|\le C\), \(b:H\to H\) measurable, \(x,z:\Omega\to H\) and \(w:\Omega\to\mathbb R\) measurable with \(|w|\le A\) almost surely (\(A\ge0\)), and assume that the joint law of \(((w,x),z)\) is the product of the law of \((w,x)\) and \(\gamma\) (that is, \(z\sim\gamma\) is independent of \((w,x)\)). Then for every \(h\in\mathbb R\) (with \(\sqrt{2h}=0\) for \(h<0\)),


\[\mathbb E\bigl[w\,\varphi\bigl(x+hb(x)+\sqrt{2h}\,z\bigr)\bigr]=\mathbb E\Bigl[w\int_H\varphi\bigl(x+hb(x)+\sqrt{2h}\,y\bigr)\,\gamma(dy)\Bigr].\]


(3) In the setting of (2), let in addition \(\varphi\in C^2\), \(b(x)\in L_2(P)\), and \(w\,(D\varphi(x)b(x)+\Delta\varphi(x))\) integrable. Let \(\varepsilon>0\), \(K\ge0\), and assume the uniform Gaussian Taylor estimate


\[\Bigl|\int\varphi(y+hv+\sqrt{2h}\,u)\,\gamma(du)-\varphi(y)-h\bigl(D\varphi(y)v+\Delta\varphi(y)\bigr)\Bigr|\le h(\varepsilon+K\sqrt h)(1+\|v\|^2)\qquad(1)\]


for all \(y,v\in H\) and \(0\le h\le1\). Then for \(0\le h\le1\),


\[\begin{gathered}
\left|\begin{aligned}
&\mathbb E\bigl[w\varphi(x+hb(x)+\sqrt{2h}z)\bigr]-\mathbb E[w\varphi(x)]\\
&\quad-h\,\mathbb E\bigl[w\bigl(D\varphi(x)b(x)+\Delta\varphi(x)\bigr)\bigr]
\end{aligned}\right|\\
\le A\,h\,(\varepsilon+K\sqrt h)\bigl(1+\mathbb E\|b(x)\|^2\bigr).
\end{gathered}\]


\textbf{Proof.} (1) \(wF\) is measurable and \(|wF|\le AC\) almost surely; a constant is integrable for a probability measure, so \(wF\) is integrable by domination.




(2) Let \(G((s,u),y)=s\,\varphi(u+hb(u)+\sqrt{2h}y)\) on \((\mathbb R\times H)\times H\), a measurable function. By the change-of-variables formula and the product form of the law,


\[\mathbb E\bigl[G((w,x),z)\bigr]=\int G\,d\bigl(\mathcal L(w,x)\otimes\gamma\bigr)=\int_{\Omega\times H}w(a)\,\varphi\bigl(x(a)+hb(x(a))+\sqrt{2h}\,y\bigr)\,(P\otimes\gamma)(da,dy),\]


because \(\mathcal L(w,x)\otimes\gamma\) is the image of \(P\otimes\gamma\) under \((a,y)\mapsto((w(a),x(a)),y)\). The last integrand is integrable by (1), applied on \(\Omega\times H\) with \(P\otimes\gamma\) (the bound \(|w(a)|\le A\) holds for \(P\otimes\gamma\)-almost every \((a,y)\), the first marginal being \(P\)). Fubini's theorem integrates first in \(y\), and the factor \(w(a)\) comes out of the inner integral.




(3) Multiply (1) at \(y=x(a)\), \(v=b(x(a))\) by \(|w(a)|\le A\) and integrate; this gives (as proved in the section on compact \(C^2\) Taylor estimates)


\[\begin{gathered}
\left|\mathbb E\left[w\left(\begin{aligned}
&\int\varphi(x+hb(x)+\sqrt{2h}u)\gamma(du)-\varphi(x)\\
&\quad-h\bigl(D\varphi(x)b(x)+\Delta\varphi(x)\bigr)
\end{aligned}\right)\right]\right|\\
\le Ah(\varepsilon+K\sqrt h)\bigl(1+\mathbb E\|b(x)\|^2\bigr).
\end{gathered}\]


The three functions \(w\int\varphi(\cdot)\,d\gamma\), \(w\varphi(x)\) and \(w(D\varphi(x)b(x)+\Delta\varphi(x))\) are integrable: the first because its Gaussian integral is the inner integral of an integrable function on \(\Omega\times H\) (Fubini), the second by (1), the third by assumption. So the expectation splits into three expectations, and by (2) the first equals \(\mathbb E[w\varphi(x+hb(x)+\sqrt{2h}z)]\). \(\square\)

\begin{leanbox}{Lean: \texttt{integrable\_\allowbreak{}bounded\_\allowbreak{}weight\_\allowbreak{}observation}, \texttt{fresh\_\allowbreak{}weighted\_\allowbreak{}gaussian\_\allowbreak{}integral}, \texttt{fresh\_\allowbreak{}weighted\_\allowbreak{}weak\_\allowbreak{}step\_\allowbreak{}of\_\allowbreak{}bound}}
\leansource{\texttt{integrable\_\allowbreak{}bounded\_\allowbreak{}weight\_\allowbreak{}observation} \quad AverageLMC/\allowbreak{}LaplacianWeakChain.lean, lines 19–27}
\begin{leancode}
\leanline{0}{theorem\ integrable\_bounded\_weight\_observation}
\leanline{4}{(P\ :\ Measure\ Ω)\ [IsProbabilityMeasure\ P]\ (w\ F\ :\ Ω\ →\ ℝ)}
\leanline{4}{(hw\ :\ Measurable\ w)\ (hF\ :\ Measurable\ F)\ (A\ C\ :\ ℝ)\ (hA\ :\ 0\ ≤\ A)}
\leanline{4}{(hwA\ :\ ∀ᵐ\ a\ ∂P,\ |w\ a|\ ≤\ A)\ (hFC\ :\ ∀\ a,\ |F\ a|\ ≤\ C)\ :}
\leanline{4}{Integrable\ (fun\ a\ ↦\ w\ a\ *\ F\ a)\ P}
\end{leancode}
\leansource{\texttt{fresh\_\allowbreak{}weighted\_\allowbreak{}gaussian\_\allowbreak{}integral} \quad AverageLMC/\allowbreak{}LaplacianWeakChain.lean, lines 32–72}
\begin{leancode}
\leanline{0}{theorem\ fresh\_weighted\_gaussian\_integral}
\leanline{4}{(P\ :\ Measure\ Ω)\ [IsProbabilityMeasure\ P]}
\leanline{4}{(φ\ :\ H\ →\ ℝ)\ (hφ\ :\ Continuous\ φ)\ (C\ :\ ℝ)\ (hφC\ :\ ∀\ x,\ |φ\ x|\ ≤\ C)}
\leanline{4}{(b\ :\ H\ →\ H)\ (hb\ :\ Measurable\ b)\ (x\ z\ :\ Ω\ →\ H)\ (w\ :\ Ω\ →\ ℝ)}
\leanline{4}{(hx\ :\ Measurable\ x)\ (hz\ :\ Measurable\ z)\ (hw\ :\ Measurable\ w)}
\leanline{4}{(A\ :\ ℝ)\ (hA\ :\ 0\ ≤\ A)\ (hwA\ :\ ∀ᵐ\ a\ ∂P,\ |w\ a|\ ≤\ A)}
\leanline{4}{(hfresh\ :\ P.map\ (fun\ a\ ↦\ ((w\ a,\ x\ a),\ z\ a))\ =}
\leanline{6}{(P.map\ (fun\ a\ ↦\ (w\ a,\ x\ a))).prod\ (stdGaussian\ H))\ (h\ :\ ℝ)\ :}
\leanline{4}{(∫\ a,\ w\ a\ *\ φ\ (x\ a\ +\ h\ •\ b\ (x\ a)\ +\ Real.sqrt\ (2\ *\ h)\ •\ z\ a)\ ∂P)\ =}
\leanline{6}{∫\ a,\ w\ a\ *\ (∫\ y\ :\ H,\ φ\ (x\ a\ +\ h\ •\ b\ (x\ a)\ +\ Real.sqrt\ (2\ *\ h)\ •\ y)}
\leanline{8}{∂stdGaussian\ H)\ ∂P}
\end{leancode}
\leansource{\texttt{fresh\_\allowbreak{}weighted\_\allowbreak{}weak\_\allowbreak{}step\_\allowbreak{}of\_\allowbreak{}bound} \quad AverageLMC/\allowbreak{}LaplacianWeakChain.lean, lines 77–128}
\begin{leancode}
\leanline{0}{theorem\ fresh\_weighted\_weak\_step\_of\_bound}
\leanline{4}{(P\ :\ Measure\ Ω)\ [IsProbabilityMeasure\ P]}
\leanline{4}{(φ\ :\ H\ →\ ℝ)\ (hφ\ :\ ContDiff\ ℝ\ 2\ φ)\ (C\ :\ ℝ)\ (hφC\ :\ ∀\ x,\ |φ\ x|\ ≤\ C)}
\leanline{4}{(b\ :\ H\ →\ H)\ (hb\ :\ Measurable\ b)\ (x\ z\ :\ Ω\ →\ H)\ (w\ :\ Ω\ →\ ℝ)}
\leanline{4}{(hx\ :\ Measurable\ x)\ (hz\ :\ Measurable\ z)\ (hw\ :\ Measurable\ w)}
\leanline{4}{(hb₂\ :\ MemLp\ (fun\ a\ ↦\ b\ (x\ a))\ 2\ P)}
\leanline{4}{(hgen\ :\ Integrable\ (fun\ a\ ↦\ w\ a\ *}
\leanline{6}{(fderiv\ ℝ\ φ\ (x\ a)\ (b\ (x\ a))\ +\ Laplacian.laplacian\ φ\ (x\ a)))\ P)}
\leanline{4}{(A\ :\ ℝ)\ (hA\ :\ 0\ ≤\ A)\ (hwA\ :\ ∀ᵐ\ a\ ∂P,\ |w\ a|\ ≤\ A)}
\leanline{4}{(hfresh\ :\ P.map\ (fun\ a\ ↦\ ((w\ a,\ x\ a),\ z\ a))\ =}
\leanline{6}{(P.map\ (fun\ a\ ↦\ (w\ a,\ x\ a))).prod\ (stdGaussian\ H))}
\leanline{4}{(ε\ :\ ℝ)\ (hε\ :\ 0\ <\ ε)\ (K\ :\ ℝ)\ (hK\ :\ 0\ ≤\ K)}
\leanline{4}{(hbound\ :\ ∀\ (x\ v\ :\ H)\ (h\ :\ ℝ),\ 0\ ≤\ h\ →\ h\ ≤\ 1\ →}
\leanline{6}{|(∫\ z\ :\ H,\ φ\ (x\ +\ h\ •\ v\ +\ Real.sqrt\ (2\ *\ h)\ •\ z)\ ∂stdGaussian\ H)\ -}
\leanline{8}{φ\ x\ -\ h\ *\ (fderiv\ ℝ\ φ\ x\ v\ +\ Laplacian.laplacian\ φ\ x)|\ ≤}
\leanline{10}{h\ *\ (ε\ +\ K\ *\ Real.sqrt\ h)\ *\ (1\ +\ ‖v‖\ \textasciicircum{}\ 2))}
\leanline{4}{(h\ :\ ℝ)\ (hh\ :\ 0\ ≤\ h)\ (hh1\ :\ h\ ≤\ 1)\ :}
\leanline{4}{|(∫\ a,\ w\ a\ *\ φ\ (x\ a\ +\ h\ •\ b\ (x\ a)\ +\ Real.sqrt\ (2\ *\ h)\ •\ z\ a)\ ∂P)\ -}
\leanline{6}{(∫\ a,\ w\ a\ *\ φ\ (x\ a)\ ∂P)\ -}
\leanline{6}{h\ *\ (∫\ a,\ w\ a\ *\ (fderiv\ ℝ\ φ\ (x\ a)\ (b\ (x\ a))\ +\ Laplacian.laplacian\ φ\ (x\ a))\ ∂P)|\ ≤}
\leanline{6}{A\ *\ h\ *\ (ε\ +\ K\ *\ Real.sqrt\ h)\ *\ (1\ +\ ∫\ a,\ ‖b\ (x\ a)‖\ \textasciicircum{}\ 2\ ∂P)}
\end{leancode}
\leanlabel{How the Lean proof follows the argument}
\begin{enumerate}


\item \texttt{integrable\_\allowbreak{}bounded\_\allowbreak{}weight\_\allowbreak{}observation} is (1): domination of \(wF\) by the integrable constant \texttt{A * C} (\texttt{integrable\_\allowbreak{}const} and \texttt{Integrable.\allowbreak{}mono'}), with the pointwise bound from \texttt{abs\_\allowbreak{}mul} and \texttt{mul\_\allowbreak{}le\_\allowbreak{}mul}.

\item \texttt{fresh\_\allowbreak{}weighted\_\allowbreak{}gaussian\_\allowbreak{}integral} is (2). The freshness hypothesis is \texttt{hfresh : P.\allowbreak{}map (fun a \ensuremath{\mapsto} ((w a, x a), z a)) = (P.\allowbreak{}map (fun a \ensuremath{\mapsto} (w a, x a))).\allowbreak{}prod (std\allowbreak{}Gaussian H)}, and \texttt{std\allowbreak{}Gaussian H} is \(\gamma\). \texttt{F} is the function \(G\); \texttt{hmap} is the image identity from \texttt{Measure.\allowbreak{}map\_\allowbreak{}prod\_\allowbreak{}map}; \texttt{heq} is the displayed chain, a \texttt{calc} with two applications of \texttt{integral\_\allowbreak{}map} and the rewrite by \texttt{hfresh}. \texttt{hi} is the integrability on the product (part (1) together with \texttt{measure\allowbreak{}Preserving\_\allowbreak{}fst} to transfer the almost-sure bound), \texttt{integral\_\allowbreak{}prod} is Fubini's theorem, and \texttt{integral\_\allowbreak{}const\_\allowbreak{}mul} extracts \(w(a)\).

\item \texttt{fresh\_\allowbreak{}weighted\_\allowbreak{}weak\_\allowbreak{}step\_\allowbreak{}of\_\allowbreak{}bound} is (3); the hypothesis \texttt{hbound} is (1), with \texttt{fderiv \ensuremath{\mathbb R} \ensuremath{\varphi} x v} for \(D\varphi(x)v\) and \texttt{Laplacian.\allowbreak{}laplacian \ensuremath{\varphi} x} for \(\Delta\varphi(x)\). \texttt{hweak} is the integrated estimate \texttt{gaussian\_\allowbreak{}weighted\_\allowbreak{}weak\_\allowbreak{}step\_\allowbreak{}of\_\allowbreak{}bound}. \texttt{hprod}, \texttt{hmean} (integrability of the Gaussian average, from \texttt{integral\_\allowbreak{}prod\_\allowbreak{}left}) and \texttt{hbase} provide the integrability, \texttt{heq} regroups the integrand, two applications of \texttt{integral\_\allowbreak{}sub} and one of \texttt{integral\_\allowbreak{}const\_\allowbreak{}mul} split the expectation, and the final rewrite uses part (2).

\end{enumerate}
\end{leanbox}

\section{\texorpdfstring{Freshness along the chain and the one-step residual}{Freshness along the chain and the one-step residual}}


\textbf{Setting.} Let \(U:H\to\mathbb R\) have an \(L\)-Lipschitz gradient, \(h\in\mathbb R\), \(X_0\) and \(W_0,W_1,\dots\) square-integrable random vectors in \(H\), and \(w\) a measurable real random variable, such that \((w,X_0)\) is independent of the sequence \((W_k)_{k\ge0}\), the \(W_k\) are mutually independent, and each \(W_k\) has law \(\gamma\). The Euler chain is \(X_{k+1}=X_k-h\nabla U(X_k)+\sqrt{2h}\,W_k\). For a function \(F:H\to\mathbb R\) write \(T(F,X)=\mathbb E[w\,F(X)]\), and let \(\mathcal L\varphi=\Delta\varphi-\langle\nabla U,\nabla\varphi\rangle\).




\textbf{Claim.} (1) For every \(k\), the joint law of \(((w,X_k),W_k)\) is the product of the law of \((w,X_k)\) and \(\gamma\).




(2) Let \(\varphi\in C^2\) with \(|\varphi|\le C\) and \(|\mathcal L\varphi|\le D\), let \(|w|\le A\) almost surely with \(A\ge0\), and let \(\varphi\) satisfy (1) of Step 1 with \(\varepsilon>0\), \(K\ge0\). Then for \(0\le h\le1\) and every \(k\),


\[\bigl|T(\varphi,X_{k+1})-T(\varphi,X_k)-h\,T(\mathcal L\varphi,X_k)\bigr|\le A\,h\,(\varepsilon+K\sqrt h)\bigl(1+\|\nabla U(X_k)\|_{L_2}^2\bigr).\]


\textbf{Proof.} (1) Consider the augmented chain \(Z_k=(w,X_k)\) with values in \(\mathbb R\times H\). It satisfies \(Z_{k+1}=\Phi(Z_k,W_k)\) with \(\Phi((s,u),y)=(s,\,u-h\nabla U(u)+\sqrt{2h}y)\), a measurable map, since the first coordinate is carried along unchanged. For a chain driven by independent noise with common law \(\gamma\) and an initial state independent of the noise, \(Z_k\) is a measurable function of \(Z_0,W_0,\dots,W_{k-1}\), which are independent of \(W_k\); hence the joint law of \((Z_k,W_k)\) is the product of the law of \(Z_k\) and \(\gamma\) (this is proved in the section on independent Euler steps). Since \(Z_k=(w,X_k)\), this is the claim.




(2) Apply (3) of Step 1 with \(x=X_k\), \(z=W_k\) and \(b=-\nabla U\). Its hypotheses hold: \(-\nabla U\) is continuous, hence measurable; the freshness of \(W_k\) is part (1); \(b(X_k)=-\nabla U(X_k)\) is square integrable because a Lipschitz map sends square-integrable vectors to square-integrable vectors (\(\|\nabla U(u)\|\le\|\nabla U(0)\|+L\|u\|\)). Since \(D\varphi(u)v=\langle\nabla\varphi(u),v\rangle\),


\[D\varphi(u)\bigl(-\nabla U(u)\bigr)+\Delta\varphi(u)=\Delta\varphi(u)-\langle\nabla U(u),\nabla\varphi(u)\rangle=\mathcal L\varphi(u),\]


and \(w\,\mathcal L\varphi(X_k)\) is integrable by (1) of Step 1, because \(\mathcal L\varphi\) is continuous and bounded by \(D\). Finally \(X_{k+1}=X_k+h(-\nabla U(X_k))+\sqrt{2h}W_k\) almost surely, so \(T(\varphi,X_{k+1})\) is the first expectation in (3) of Step 1, and \(\mathbb E\|b(X_k)\|^2=\|\nabla U(X_k)\|_{L_2}^2\). \(\square\)

\begin{leanbox}{Lean: \texttt{hilbert\_\allowbreak{}lpChain\_\allowbreak{}weighted\_\allowbreak{}fresh}, \texttt{hilbert\_\allowbreak{}lpChain\_\allowbreak{}weak\_\allowbreak{}step\_\allowbreak{}of\_\allowbreak{}bound}}
\leansource{\texttt{hilbert\_\allowbreak{}lpChain\_\allowbreak{}weighted\_\allowbreak{}fresh} \quad AverageLMC/\allowbreak{}LaplacianWeakChain.lean, lines 146–178}
\begin{leancode}
\leanline{0}{theorem\ hilbert\_lpChain\_weighted\_fresh}
\leanline{4}{(P\ :\ Measure\ Ω)\ [IsProbabilityMeasure\ P]}
\leanline{4}{(U\ :\ H\ →\ ℝ)\ (h\ :\ ℝ)\ (L\ :\ ℝ≥0)\ (hLip\ :\ LipschitzWith\ L\ (gradient\ U))}
\leanline{4}{(initial\ :\ Lp\ H\ 2\ P)\ (W\ :\ ℕ\ →\ Lp\ H\ 2\ P)\ (w\ :\ Ω\ →\ ℝ)\ (hw\ :\ Measurable\ w)}
\leanline{4}{(hind\ :\ IndepFun\ (fun\ a\ ↦\ (w\ a,\ initial\ a))\ (fun\ a\ k\ ↦\ W\ k\ a)\ P)}
\leanline{4}{(hiid\ :\ iIndepFun\ (fun\ k\ a\ ↦\ W\ k\ a)\ P)}
\leanline{4}{(hlaw\ :\ ∀\ k,\ P.map\ (W\ k)\ =\ stdGaussian\ H)\ (k\ :\ ℕ)\ :}
\leanline{4}{P.map\ (fun\ a\ ↦\ ((w\ a,\ ULA.Hilbert.lpChain\ P\ U\ h\ hLip\ initial\ W\ k\ a),\ W\ k\ a))\ =}
\leanline{6}{(P.map\ (fun\ a\ ↦\ (w\ a,\ ULA.Hilbert.lpChain\ P\ U\ h\ hLip\ initial\ W\ k\ a))).prod}
\leanline{8}{(stdGaussian\ H)}
\end{leancode}
\leansource{\texttt{hilbert\_\allowbreak{}lpChain\_\allowbreak{}weak\_\allowbreak{}step\_\allowbreak{}of\_\allowbreak{}bound} \quad AverageLMC/\allowbreak{}LaplacianWeakChain.lean, lines 182–239}
\begin{leancode}
\leanline{0}{theorem\ hilbert\_lpChain\_weak\_step\_of\_bound}
\leanline{4}{(P\ :\ Measure\ Ω)\ [IsProbabilityMeasure\ P]}
\leanline{4}{(U\ :\ H\ →\ ℝ)\ (h\ :\ ℝ)\ (L\ :\ ℝ≥0)\ (hLip\ :\ LipschitzWith\ L\ (gradient\ U))}
\leanline{4}{(initial\ :\ Lp\ H\ 2\ P)\ (W\ :\ ℕ\ →\ Lp\ H\ 2\ P)\ (w\ :\ Ω\ →\ ℝ)\ (hw\ :\ Measurable\ w)}
\leanline{4}{(hind\ :\ IndepFun\ (fun\ a\ ↦\ (w\ a,\ initial\ a))\ (fun\ a\ k\ ↦\ W\ k\ a)\ P)}
\leanline{4}{(hiid\ :\ iIndepFun\ (fun\ k\ a\ ↦\ W\ k\ a)\ P)}
\leanline{4}{(hlaw\ :\ ∀\ k,\ P.map\ (W\ k)\ =\ stdGaussian\ H)}
\leanline{4}{(φ\ :\ H\ →\ ℝ)\ (hφ\ :\ ContDiff\ ℝ\ 2\ φ)}
\leanline{4}{(C\ D\ :\ ℝ)\ (hφC\ :\ ∀\ x,\ |φ\ x|\ ≤\ C)\ (hGD\ :\ ∀\ x,\ |ULA.langevinGenerator\ U\ φ\ x|\ ≤\ D)}
\leanline{4}{(A\ :\ ℝ)\ (hA\ :\ 0\ ≤\ A)\ (hwA\ :\ ∀ᵐ\ a\ ∂P,\ |w\ a|\ ≤\ A)}
\leanline{4}{(ε\ :\ ℝ)\ (hε\ :\ 0\ <\ ε)\ (K\ :\ ℝ)\ (hK\ :\ 0\ ≤\ K)}
\leanline{4}{(hbound\ :\ ∀\ (x\ v\ :\ H)\ (h\ :\ ℝ),\ 0\ ≤\ h\ →\ h\ ≤\ 1\ →}
\leanline{6}{|(∫\ z\ :\ H,\ φ\ (x\ +\ h\ •\ v\ +\ Real.sqrt\ (2\ *\ h)\ •\ z)\ ∂stdGaussian\ H)\ -}
\leanline{8}{φ\ x\ -\ h\ *\ (fderiv\ ℝ\ φ\ x\ v\ +\ Laplacian.laplacian\ φ\ x)|\ ≤}
\leanline{10}{h\ *\ (ε\ +\ K\ *\ Real.sqrt\ h)\ *\ (1\ +\ ‖v‖\ \textasciicircum{}\ 2))}
\leanline{4}{(hh\ :\ 0\ ≤\ h)\ (hh1\ :\ h\ ≤\ 1)\ (k\ :\ ℕ)\ :}
\leanline{4}{|lpWeightedTest\ P\ w\ φ\ (ULA.Hilbert.lpChain\ P\ U\ h\ hLip\ initial\ W\ (k\ +\ 1))\ -}
\leanline{6}{lpWeightedTest\ P\ w\ φ\ (ULA.Hilbert.lpChain\ P\ U\ h\ hLip\ initial\ W\ k)\ -}
\leanline{6}{h\ *\ lpWeightedTest\ P\ w\ (ULA.langevinGenerator\ U\ φ)\ (ULA.Hilbert.lpChain\ P\ U\ h\ hLip\ initial\ W\ k)|\ ≤}
\leanline{6}{A\ *\ h\ *\ (ε\ +\ K\ *\ Real.sqrt\ h)\ *}
\leanline{8}{(1\ +\ ‖ULA.lpDrift\ P\ (gradient\ U)\ hLip\ (ULA.Hilbert.lpChain\ P\ U\ h\ hLip\ initial\ W\ k)‖\ \textasciicircum{}\ 2)}
\end{leancode}
\leanlabel{How the Lean proof follows the argument}
The chain is \texttt{ULA.\allowbreak{}Hilbert.\allowbreak{}lp\allowbreak{}Chain P U h h\allowbreak{}Lip initial W k}, a sequence in the Hilbert space \texttt{Lp H 2 P} of square-integrable classes, defined by recursion with the update \texttt{ULA.\allowbreak{}lp\allowbreak{}Ula\allowbreak{}Update}; \texttt{ULA.\allowbreak{}hilbert\_\allowbreak{}lp\allowbreak{}Chain\_\allowbreak{}coe} states that it agrees almost surely with the pointwise recursion \texttt{ULA.\allowbreak{}noise\allowbreak{}Driven\allowbreak{}Chain (ULA.\allowbreak{}Hilbert.\allowbreak{}pointwise\allowbreak{}Update U h) (initial \ensuremath{\omega}) (fun k \ensuremath{\mapsto} W k \ensuremath{\omega}) n}. Independence is \texttt{hind : Indep\allowbreak{}Fun (fun a \ensuremath{\mapsto} (w a, initial a)) (fun a k \ensuremath{\mapsto} W k a) P} and \texttt{hiid : i\allowbreak{}Indep\allowbreak{}Fun (fun k a \ensuremath{\mapsto} W k a) P}, and \texttt{hlaw} gives the Gaussian laws.



\begin{enumerate}


\item \texttt{hilbert\_\allowbreak{}lp\allowbreak{}Chain\_\allowbreak{}weighted\_\allowbreak{}fresh} is (1). \texttt{Fp} is the augmented update \(\Phi\), measurable by \texttt{ULA.\allowbreak{}Hilbert.\allowbreak{}pointwise\allowbreak{}Update\_\allowbreak{}measurable}; \texttt{hj} is \texttt{noise\allowbreak{}Driven\allowbreak{}Chain\_\allowbreak{}joint\_\allowbreak{}fresh\_\allowbreak{}of\_\allowbreak{}iid} for the augmented chain, and \texttt{noise\allowbreak{}Driven\allowbreak{}Chain\_\allowbreak{}passive} rewrites the augmented chain as \((w,X_k)\). \texttt{hm\ensuremath{{}_{1}}} and \texttt{hm\ensuremath{{}_{2}}} replace the pointwise chain by the \(L^2\) chain inside both laws, using \texttt{Measure.\allowbreak{}map\_\allowbreak{}congr} and the almost-sure identity \texttt{hc}.

\item \texttt{hilbert\_\allowbreak{}lp\allowbreak{}Chain\_\allowbreak{}weak\_\allowbreak{}step\_\allowbreak{}of\_\allowbreak{}bound} is (2), with \texttt{lp\allowbreak{}Weighted\allowbreak{}Test P w \ensuremath{\varphi} X} for \(T(\varphi,X)\) and \texttt{ULA.\allowbreak{}langevin\allowbreak{}Generator U \ensuremath{\varphi}} for \(\mathcal L\varphi\). \texttt{h\allowbreak{}G} is the generator identity, from \texttt{ULA.\allowbreak{}langevin\allowbreak{}Generator\_\allowbreak{}eq\_\allowbreak{}laplacian\_\allowbreak{}sub\_\allowbreak{}inner} and \texttt{inner\_\allowbreak{}gradient\_\allowbreak{}left}; \texttt{h\allowbreak{}Gc} is continuity of the generator; \texttt{hgen} the integrability of \(w\mathcal L\varphi(X_k)\). \texttt{hweak} is Step 1 (3) applied with freshness from part (1) and square integrability from \texttt{ULA.\allowbreak{}lp\allowbreak{}Drift\_\allowbreak{}mem\allowbreak{}Lp}. \texttt{hnext} identifies \(T(\varphi,X_{k+1})\) through \texttt{ULA.\allowbreak{}lp\allowbreak{}Ula\allowbreak{}Update\_\allowbreak{}coe}, and \texttt{hnorm} computes the squared norm of the \(L^2\) drift \texttt{ULA.\allowbreak{}lp\allowbreak{}Drift P (gradient U) h\allowbreak{}Lip X} as \(\mathbb E\|\nabla U(X_k)\|^2\) via \texttt{L2.\allowbreak{}inner\_\allowbreak{}def}.

\end{enumerate}
\end{leanbox}


Earlier results used here: \cite[Section~3.13.11]{LangevinLean2026ULASC}; \cite[Section~3.13.12]{LangevinLean2026ULASC}; \cite[Section~3.13.2]{LangevinLean2026ULASC}; \cite[Section~3.13.3]{LangevinLean2026ULASC}; \cite[Section~3.13.4]{LangevinLean2026ULASC}; \cite[Section~3.19.1]{LangevinLean2026ULASC}; \cite[Section~3.20.1]{LangevinLean2026ULASC}; \cite[Section~3.23.1]{LangevinLean2026ULASC}; \cite[Section~3.23.2]{LangevinLean2026ULASC}; \cite[Section~3.42.1]{LangevinLean2026ULASC}; \cite[Section~3.46.14]{LangevinLean2026ULASC}; \cite[Section~3.46.2]{LangevinLean2026ULASC}; \cite[Section~3.46.8]{LangevinLean2026ULASC}.

\section{\texorpdfstring{Summing the residuals over a time horizon}{Summing the residuals over a time horizon}}


\textbf{Claim.} (1) Let \((a_k)\) and \((c_k)\) be real sequences, \(h\in\mathbb R\), \(E\ge0\) and \(n\in\mathbb N\). If \(|a_{k+1}-a_k-hc_k|\le E\) for every \(k<n\), then


\[\Bigl|a_n-a_0-\sum_{k=0}^{n-1}h\,c_k\Bigr|\le nE.\]


(2) In the setting of Step 2 (2), let \(R\ge0\) with \(\|\nabla U(X_k)\|_{L_2}\le R\) for every \(k\le n\). Then


\[\Bigl|T(\varphi,X_n)-T(\varphi,X_0)-\sum_{k=0}^{n-1}h\,T(\mathcal L\varphi,X_k)\Bigr|\le A\,(nh)\,(\varepsilon+K\sqrt h)\,(1+R^2).\]


\textbf{Proof.} (1) The sum telescopes: \(\sum_{k=0}^{n-1}(a_{k+1}-a_k-hc_k)=a_n-a_0-\sum_{k=0}^{n-1}hc_k\). By the triangle inequality for finite sums,


\[\Bigl|a_n-a_0-\sum_{k<n}hc_k\Bigr|\le\sum_{k<n}|a_{k+1}-a_k-hc_k|\le\sum_{k<n}E=nE.\]


(2) For \(k<n\) we have \(\|\nabla U(X_k)\|_{L_2}^2\le R^2\), so by Step 2 (2) each local residual is at most \(E=Ah(\varepsilon+K\sqrt h)(1+R^2)\), which is nonnegative. Part (1) with \(a_k=T(\varphi,X_k)\) and \(c_k=T(\mathcal L\varphi,X_k)\) gives the bound \(nE=A(nh)(\varepsilon+K\sqrt h)(1+R^2)\). \(\square\)




In the Dynkin identity, \(nh\) is the fixed time horizon \(t\), and the right-hand side tends to \(At\,\varepsilon(1+R^2)\) as \(h\to0\); since \(\varepsilon>0\) is arbitrary, the Euler sums approximate \(\int_0^tT(\mathcal L\varphi,\cdot)\,ds\).

\begin{leanbox}{Lean: \texttt{sum\_\allowbreak{}weighted\_\allowbreak{}weak\_\allowbreak{}residual}, \texttt{hilbert\_\allowbreak{}lpChain\_\allowbreak{}weak\_\allowbreak{}sum\_\allowbreak{}of\_\allowbreak{}bound}}
\leansource{\texttt{sum\_\allowbreak{}weighted\_\allowbreak{}weak\_\allowbreak{}residual} \quad AverageLMC/\allowbreak{}LaplacianWeakChain.lean, lines 132–142}
\begin{leancode}
\leanline{0}{theorem\ sum\_weighted\_weak\_residual\ (a\ c\ :\ ℕ\ →\ ℝ)\ (h\ E\ :\ ℝ)\ (n\ :\ ℕ)}
\leanline{4}{(\_hE\ :\ 0\ ≤\ E)\ (hlocal\ :\ ∀\ k\ <\ n,\ |a\ (k\ +\ 1)\ -\ a\ k\ -\ h\ *\ c\ k|\ ≤\ E)\ :}
\leanline{4}{|a\ n\ -\ a\ 0\ -\ ∑\ k\ ∈\ Finset.range\ n,\ h\ *\ c\ k|\ ≤\ (n\ :\ ℝ)\ *\ E}
\end{leancode}
\leansource{\texttt{hilbert\_\allowbreak{}lpChain\_\allowbreak{}weak\_\allowbreak{}sum\_\allowbreak{}of\_\allowbreak{}bound} \quad AverageLMC/\allowbreak{}LaplacianWeakChain.lean, lines 244–284}
\begin{leancode}
\leanline{0}{theorem\ hilbert\_lpChain\_weak\_sum\_of\_bound}
\leanline{4}{(P\ :\ Measure\ Ω)\ [IsProbabilityMeasure\ P]}
\leanline{4}{(U\ :\ H\ →\ ℝ)\ (h\ :\ ℝ)\ (L\ :\ ℝ≥0)\ (hLip\ :\ LipschitzWith\ L\ (gradient\ U))}
\leanline{4}{(initial\ :\ Lp\ H\ 2\ P)\ (W\ :\ ℕ\ →\ Lp\ H\ 2\ P)\ (w\ :\ Ω\ →\ ℝ)\ (hw\ :\ Measurable\ w)}
\leanline{4}{(hind\ :\ IndepFun\ (fun\ a\ ↦\ (w\ a,\ initial\ a))\ (fun\ a\ k\ ↦\ W\ k\ a)\ P)}
\leanline{4}{(hiid\ :\ iIndepFun\ (fun\ k\ a\ ↦\ W\ k\ a)\ P)}
\leanline{4}{(hlaw\ :\ ∀\ k,\ P.map\ (W\ k)\ =\ stdGaussian\ H)}
\leanline{4}{(φ\ :\ H\ →\ ℝ)\ (hφ\ :\ ContDiff\ ℝ\ 2\ φ)}
\leanline{4}{(C\ D\ :\ ℝ)\ (hφC\ :\ ∀\ x,\ |φ\ x|\ ≤\ C)\ (hGD\ :\ ∀\ x,\ |ULA.langevinGenerator\ U\ φ\ x|\ ≤\ D)}
\leanline{4}{(A\ :\ ℝ)\ (hA\ :\ 0\ ≤\ A)\ (hwA\ :\ ∀ᵐ\ a\ ∂P,\ |w\ a|\ ≤\ A)}
\leanline{4}{(ε\ :\ ℝ)\ (hε\ :\ 0\ <\ ε)\ (K\ :\ ℝ)\ (hK\ :\ 0\ ≤\ K)}
\leanline{4}{(hbound\ :\ ∀\ (x\ v\ :\ H)\ (h\ :\ ℝ),\ 0\ ≤\ h\ →\ h\ ≤\ 1\ →}
\leanline{6}{|(∫\ z\ :\ H,\ φ\ (x\ +\ h\ •\ v\ +\ Real.sqrt\ (2\ *\ h)\ •\ z)\ ∂stdGaussian\ H)\ -}
\leanline{8}{φ\ x\ -\ h\ *\ (fderiv\ ℝ\ φ\ x\ v\ +\ Laplacian.laplacian\ φ\ x)|\ ≤}
\leanline{10}{h\ *\ (ε\ +\ K\ *\ Real.sqrt\ h)\ *\ (1\ +\ ‖v‖\ \textasciicircum{}\ 2))}
\leanline{4}{(hh\ :\ 0\ ≤\ h)\ (hh1\ :\ h\ ≤\ 1)\ (n\ :\ ℕ)\ (R\ :\ ℝ)\ (hR\ :\ 0\ ≤\ R)}
\leanline{4}{(hboundR\ :\ ∀\ k\ ≤\ n,}
\leanline{6}{‖ULA.lpDrift\ P\ (gradient\ U)\ hLip\ (ULA.Hilbert.lpChain\ P\ U\ h\ hLip\ initial\ W\ k)‖\ ≤\ R)\ :}
\leanline{4}{|lpWeightedTest\ P\ w\ φ\ (ULA.Hilbert.lpChain\ P\ U\ h\ hLip\ initial\ W\ n)\ -}
\leanline{6}{lpWeightedTest\ P\ w\ φ\ initial\ -\ ∑\ k\ ∈\ Finset.range\ n,}
\leanline{8}{h\ *\ lpWeightedTest\ P\ w\ (ULA.langevinGenerator\ U\ φ)}
\leanline{10}{(ULA.Hilbert.lpChain\ P\ U\ h\ hLip\ initial\ W\ k)|\ ≤}
\leanline{6}{A\ *\ ((n\ :\ ℝ)\ *\ h)\ *\ (ε\ +\ K\ *\ Real.sqrt\ h)\ *\ (1\ +\ R\ \textasciicircum{}\ 2)}
\end{leancode}
\leanlabel{How the Lean proof follows the argument}
\begin{enumerate}


\item \texttt{sum\_\allowbreak{}weighted\_\allowbreak{}weak\_\allowbreak{}residual} has the hypothesis \texttt{hlocal : \ensuremath{\forall} k < n, |a (k + 1) -\allowbreak{} a k -\allowbreak{} h * c k| \ensuremath{\le} E} and an unused sign condition \texttt{\_\allowbreak{}h\allowbreak{}E}. \texttt{ht} is the telescoping identity (\texttt{Finset.\allowbreak{}sum\_\allowbreak{}sub\_\allowbreak{}distrib} and \texttt{Finset.\allowbreak{}sum\_\allowbreak{}range\_\allowbreak{}sub}), and the \texttt{calc} applies \texttt{Finset.\allowbreak{}abs\_\allowbreak{}sum\_\allowbreak{}le\_\allowbreak{}sum\_\allowbreak{}abs} and \texttt{Finset.\allowbreak{}sum\_\allowbreak{}le\_\allowbreak{}sum}.

\item \texttt{hilbert\_\allowbreak{}lp\allowbreak{}Chain\_\allowbreak{}weak\_\allowbreak{}sum\_\allowbreak{}of\_\allowbreak{}bound} adds \texttt{h\allowbreak{}R : 0 \ensuremath{\le} R} and \texttt{hbound\allowbreak{}R}, the bound on the \(L^2\) drift norms for \(k\le n\). \texttt{hlocal} combines \texttt{hilbert\_\allowbreak{}lp\allowbreak{}Chain\_\allowbreak{}weak\_\allowbreak{}step\_\allowbreak{}of\_\allowbreak{}bound} with the squared bound \texttt{hs} (\texttt{sq\_\allowbreak{}le\_\allowbreak{}sq\ensuremath{{}_{0}}}) and \texttt{mul\_\allowbreak{}le\_\allowbreak{}mul\_\allowbreak{}of\_\allowbreak{}nonneg\_\allowbreak{}left}; part (1) is applied as \texttt{hb}, and the final \texttt{simpa} rearranges \(n\cdot Ah(\cdots)\) as \(A(nh)(\cdots)\).

\end{enumerate}
\end{leanbox}


Earlier results used here: \cite[Section~3.13.2]{LangevinLean2026ULASC}; \cite[Section~3.19.1]{LangevinLean2026ULASC}; \cite[Section~3.46.2]{LangevinLean2026ULASC}.

\begin{thebibliography}{ULA-LSI}
\bibitem[ULA-SC]{LangevinLean2026ULASC}
LangevinLean Project. \emph{Unadjusted Langevin under Strong Convexity}. LL-01, 2026. Draft report, content revision d50c0f64815b. \url{https://odakyildiz.github.io/LangevinLean/docs/pdf/UnadjustedLangevin/FullProof.pdf}.
\end{thebibliography}

\end{document}
